%! AP Statistics — Unit 2, Lesson 2.5 — Warm-Up (Answer Key)
\documentclass[10pt]{article}
\usepackage{apstats-article}
\usepackage{apstats-key}

\begin{document}

\pageheader{Unit 2, Lesson 2.5}{Warm-Up --- Answer Key}
\namedateperiod

\vspace{0.2in}

{\large\bfseries\color{navy} Part 1 \quad DOFS Review (Lesson 2.4)}

\vspace{0.1in}

\noindent For each scenario, write a complete DOFS description. Use the words
\textit{strong}, \textit{moderate}, or \textit{weak} and \textit{positive},
\textit{negative}, or \textit{none}.

\vspace{0.08in}

\begin{enumerate}[leftmargin=1.6em, itemsep=0.18in]

  \item A scatterplot of hours of sleep ($x$) vs.\ reaction time in milliseconds ($y$)
        for 20 college students shows points that closely follow a line sloping downward,
        with no obvious outliers.\\[8pt]
        \ansline{There is a strong negative linear association between hours of sleep and
        reaction time (ms) for these 20 college students. As sleep increases, reaction time
        tends to decrease. There are no obvious outliers.}

  \item A scatterplot of shoe size ($x$) vs.\ GPA ($y$) for 30 students shows points
        scattered randomly with no discernible pattern.\\[8pt]
        \ansline{There is no association between shoe size and GPA for these 30 students.
        Knowing a student's shoe size gives no information about their GPA.}

\end{enumerate}

\vspace{0.22in}

{\large\bfseries\color{navy} Part 2 \quad Estimating the Correlation}

\vspace{0.1in}

\noindent The \textbf{correlation coefficient} $r$ is a single number between $-1$ and $1$
that captures the \emph{direction} and \emph{strength} of a linear association.

\vspace{0.1in}

\noindent Match each description to the $r$ value you think fits best.

\vspace{0.1in}

\begin{center}
\renewcommand{\arraystretch}{1.6}
\begin{tabularx}{0.85\linewidth}{@{} l X c @{}}
\rowcolor{sky}
\textbf{Description} & \textbf{What you expect} & \textbf{$r$ value} \\
\midrule
Strong positive linear association  & points tightly around a rising line  & \ans{$0.92$} \\
No linear association               & randomly scattered cloud              & \ans{$0$} \\
Moderate negative linear association & loosely around a falling line       & \ans{$-0.55$} \\
Perfect positive linear association  & every point exactly on a line       & \ans{$1.0$} \\
\end{tabularx}
\end{center}

\vspace{0.1in}

\noindent \textbf{Values to choose from:} $\quad r = 1.0 \qquad r = 0.92 \qquad r = 0 \qquad r = -0.55$

\vspace{0.18in}

\noindent Based on your answers, predict: in Lesson 2.4's sleep vs.\ reaction-time
scatterplot (Part 1, \#1), would $r$ be closer to $-1$, $0$, or $+1$? Why?\\[8pt]
\ansline{$r$ would be closer to $-1$ because the association is strong and negative ---
as sleep increases, reaction time decreases, so the two variables move in opposite directions.}

\end{document}
